\chapter*{Tóm tắt}
\addcontentsline{toc}{chapter}{Tóm tắt}

\textbf{Tóm tắt:} Web Application Firewall (WAF) là một lớp phòng thủ quan trọng giúp phát hiện và ngăn chặn các request độc hại trước khi chúng tới ứng dụng web. Trong các WAF mã nguồn mở, ModSecurity kết hợp với OWASP Core Rule Set (CRS) là một lựa chọn phổ biến nhờ khả năng phát hiện nhiều nhóm tấn công như Cross-Site Scripting (XSS), SQL Injection (SQLi), Local File Inclusion, Remote Code Execution và các bất thường trong giao thức HTTP. Tuy nhiên, khi CRS được cấu hình ở mức nghiêm ngặt cao, đặc biệt là Paranoia Level 4, số lượng false positives có thể tăng mạnh, dẫn đến việc chặn nhầm traffic hợp lệ và làm giảm khả năng triển khai thực tế của WAF. Bên cạnh đó, các kỹ thuật tấn công hiện đại thường sử dụng encoding nhiều lớp, obfuscation, biến thể cú pháp JavaScript/SQL hoặc payload được thiết kế để né tránh regex, khiến các rule dựa trên liệt kê cố định dễ bị bỏ sót.

Khóa luận này nghiên cứu bài toán tối ưu hóa bộ luật WAF ModSecurity cho phát hiện XSS và SQL Injection bằng phương pháp \textit{Generalized Rules}. Thay vì chỉ liệt kê các payload hoặc từ khóa cụ thể, khóa luận đề xuất cách tiếp cận dựa trên ngữ cảnh và cấu trúc hành vi tấn công. Đối với XSS, hệ thống tập trung vào \textit{execution context}, chẳng hạn thẻ HTML có khả năng thực thi, event handler đi kèm hàm nguy hiểm, protocol URI như \path|javascript:|, DOM sink và các biến thể obfuscation. Đối với SQLi, hệ thống tập trung vào các \textit{semantic patterns} như \path|UNION SELECT|, stacked queries, boolean-based injection, time-based blind SQLi, error-based extraction, out-of-band SQLi, schema enumeration và các kỹ thuật comment/hex encoding.

Giải pháp được hiện thực dưới dạng hai bộ custom rules v2 tích hợp với cơ chế Anomaly Scoring của OWASP CRS. Bộ XSS rules v2 sử dụng dải ID \path|9942xxx| và được tổ chức theo kiến trúc 5 lớp gồm Core XSS, Structural Injection, Encoded/Obfuscated XSS, DOM Context và Advanced Heuristic, kèm một lớp xử lý Base64 riêng. Bộ SQLi rules v2 sử dụng dải ID \path|9943xxx|, phát hiện các nhóm SQLi cổ điển, blind SQLi, error-based/OOB SQLi, advanced encoding và một số heuristic liên quan tới NoSQL hoặc OS-command-via-SQL. Các rule áp dụng chuỗi transform như \path|t:urlDecodeUni|, \path|t:htmlEntityDecode|, \path|t:jsDecode|, \path|t:removeNulls| và \path|t:compressWhitespace| nhằm chuẩn hóa payload trước khi so khớp. Ngoài ra, khóa luận xây dựng cấu hình tuning exclusions để vô hiệu hóa có kiểm soát các CRS rules gây false positive cao hoặc trùng lặp với custom rules v2.

Để đánh giá hiệu quả, khóa luận sử dụng benchmark framework tự động gửi payload tới WAF, phân loại kết quả thành true positives, false negatives, true negatives và false positives, sau đó tính các chỉ số Recall, Precision, False Positive Rate, Accuracy và F1-score. Các tập dữ liệu được sử dụng bao gồm GoTestWAF, Mereani XSS, XSSed-DMOZ và các nguồn dữ liệu HTTP/XSS/SQLi đã được chuẩn hóa. Kết quả thực nghiệm cho thấy hướng generalized rules đặc biệt hiệu quả trong việc giảm false positives trên các tập dữ liệu mixed/benign. Trên Mereani clean, rule gốc đạt recall cao nhưng có FPR khoảng 99.78\%, trong khi custom rules v2 giữ recall khoảng 98.20\% và giảm FPR xuống khoảng 1.05\%. Trên XSSed-DMOZ clean, custom rules v2 đạt recall khoảng 99.71\%, FPR khoảng 0.23\% và accuracy khoảng 99.74\%, trong khi rule gốc có FPR khoảng 99.52\%. Các kết quả này cho thấy việc chỉ tối đa hóa detection rate là chưa đủ; với WAF thực tế, giảm false positives là điều kiện quan trọng để hệ thống có thể được triển khai và vận hành ổn định.

Khóa luận cũng thảo luận minh bạch các giới hạn và đánh đổi của phương pháp. Việc tuning bằng \path|SecRuleRemoveById| giúp giảm đáng kể false positives nhưng cần được phân tích rủi ro, đặc biệt với các CRS rules ngoài phạm vi chính của đề tài như RCE/Shell Injection hoặc các heuristic tổng quát như libinjection SQLi. Hệ thống không tuyên bố thay thế hoàn toàn OWASP CRS hoặc phát hiện 100\% mọi biến thể tấn công mới. Thay vào đó, đóng góp chính của khóa luận là một phương pháp thiết kế rule có cấu trúc, dễ giải thích, có khả năng giảm false positives đáng kể trong phạm vi XSS/SQLi, đồng thời vẫn duy trì recall cao trên các benchmark đã kiểm thử. So với các hướng AI-based detection, WAF rule-based sau tuning có ưu điểm về độ trễ thấp, khả năng giải thích rõ ràng và khả năng cập nhật nhanh bằng cách bổ sung rule mới, phù hợp với yêu cầu triển khai inline trong hệ thống web thực tế.

\vspace{8pt}
\noindent \textit{\textbf{Từ khóa}:} Web Application Firewall, ModSecurity, OWASP CRS, Cross-Site Scripting, SQL Injection, Generalized Rules, Anomaly Scoring, False Positive Reduction.
